% Part: first-order-logic
% Chapter: model-theory
% Section: overspill

\documentclass[../../../include/open-logic-section]{subfiles}

\begin{document}

\olfileid{mod}{bas}{ove}
\olsection{له بريده اوښتنه}

\begin{thm}
\ollabel{overspill} که د جملو يو سټ~$\Gamma$ د هرې کچې متناهي
مدلونه ولري، نو يو نامتناهي مدل هم لري۔
\end{thm}

\begin{proof}
د~$\Gamma$ ژبه د شمېر وړ نويو ثابتونو $c_0$، $c_1$، \dots په
ورزياتولو وغځوئ او سټ $\Gamma \cup \{c_i \neq c_j : i \neq j\}$
وګورئ۔ دا ويل چې $\Gamma$ د هرې کچې متناهي مدلونه لري، معنا ئې دا ده
چې د هر $m >0$ دپاره داسې $n\ge m$ شته چې $\Gamma$ د شمېر~$n$ يو مدل
لري۔ له دې پايله کېږي چې $\Gamma \cup \{c_i \neq c_j : i \neq j\}$
په متناهي ډول د صدق وړ دے۔ د متناهي شاهد قضيې له مخې، $\Gamma \cup
\{c_i \neq c_j : i \neq j\}$ يو مدل~$\Struct M$ لري چې د تعريف ساحه
ئې بايد نامتناهي وي، ځکه ټولې نابرابرۍ $c_i \neq c_j$ پوره کوي۔
\end{proof}

\begin{prop}
\ollabel{inf-not-fo}
د لومړۍ درجې د هېڅ ژبې داسې جمله~$!A$ نشته چې په يوه
!!a{structure}~$\Struct M$ کښې هله او يوازې هله رښتونې وي چې د دغه
!!{structure} د تعريف ساحه~$\Domain{M}$ نامتناهي وي۔
\end{prop}

\begin{proof}
که داسې~$!A$ موجوده وے، نفي~$\lnot !A$ به ئې په ټولو او يوازې په
متناهي !!{structure}s کښې رښتونې وه؛ نو د هرې کچې متناهي مدلونه به
ئې لرل خو نامتناهي مدل به ئې نۀ درلود، چې دا له~\olref{overspill}
سره په ټکر کښې ده۔
\end{proof}

\end{document}
